\documentclass[final]{article}

% if you need to pass options to natbib, use, e.g.:
%     \PassOptionsToPackage{numbers, compress}{natbib}
% before loading neurips_2025

% The authors should use one of these tracks.
% Before accepting by the NeurIPS conference, select one of the options below.
% 0. "default" for submission
 \usepackage{neurips_2025}
% the "default" option is equal to the "main" option, which is used for the Main Track with double-blind reviewing.
% 1. "main" option is used for the Main Track
%  \usepackage[main]{neurips_2025}
% 2. "position" option is used for the Position Paper Track
%  \usepackage[position]{neurips_2025}
% 3. "dandb" option is used for the Datasets & Benchmarks Track
 % \usepackage[dandb]{neurips_2025}
% 4. "creativeai" option is used for the Creative AI Track
%  \usepackage[creativeai]{neurips_2025}
% 5. "sglblindworkshop" option is used for the Workshop with single-blind reviewing
 % \usepackage[sglblindworkshop]{neurips_2025}
% 6. "dblblindworkshop" option is used for the Workshop with double-blind reviewing
%  \usepackage[dblblindworkshop]{neurips_2025}

% After being accepted, the authors should add "final" behind the track to compile a camera-ready version.
% 1. Main Track
 % \usepackage[main, final]{neurips_2025}
% 2. Position Paper Track
%  \usepackage[position, final]{neurips_2025}
% 3. Datasets & Benchmarks Track
 % \usepackage[dandb, final]{neurips_2025}
% 4. Creative AI Track
%  \usepackage[creativeai, final]{neurips_2025}
% 5. Workshop with single-blind reviewing
%  \usepackage[sglblindworkshop, final]{neurips_2025}
% 6. Workshop with double-blind reviewing
%  \usepackage[dblblindworkshop, final]{neurips_2025}
% Note. For the workshop paper template, both \title{} and \workshoptitle{} are required, with the former indicating the paper title shown in the title and the latter indicating the workshop title displayed in the footnote.
% For workshops (5., 6.), the authors should add the name of the workshop, "\workshoptitle" command is used to set the workshop title.
% \workshoptitle{WORKSHOP TITLE}

% "preprint" option is used for arXiv or other preprint submissions
 % \usepackage[preprint]{neurips_2025}

% to avoid loading the natbib package, add option nonatbib:
%    \usepackage[nonatbib]{neurips_2025}

\usepackage[utf8]{inputenc} % allow utf-8 input
\usepackage[T1]{fontenc}    % use 8-bit T1 fonts
\usepackage{hyperref}       % hyperlinks
\usepackage{url}            % simple URL typesetting
\usepackage{booktabs}       % professional-quality tables
\usepackage{amsfonts}       % blackboard math symbols
\usepackage{nicefrac}       % compact symbols for 1/2, etc.
\usepackage{microtype}      % microtypography
\usepackage{xcolor}         % colors

% Note. For the workshop paper template, both \title{} and \workshoptitle{} are required, with the former indicating the paper title shown in the title and the latter indicating the workshop title displayed in the footnote. 
\title{Rapport \\
        Projet avancé en apprentissage automatique \\
        IFT - 3710}


% The \author macro works with any number of authors. There are two commands
% used to separate the names and addresses of multiple authors: \And and \AND.
%
% Using \And between authors leaves it to LaTeX to determine where to break the
% lines. Using \AND forces a line break at that point. So, if LaTeX puts 3 of 4
% authors names on the first line, and the last on the second line, try using
% \AND instead of \And before the third author name.


\author{
    Cédric Guévremont \\
    Matricule : 20266532 \\
    \texttt{cedric.guevremont@umontreal.ca}
    \And
    Rima Boujenane \\
     Matricule : 20235550 \\
    \texttt{rima.boujenane@umontreal.ca}
    \And
    Abdelghafour Rahmouni \\
    Matricule : 20246224 \\
    \texttt{abdel.ghafour.rahmouni@umontreal.ca}
}


\begin{document}


\maketitle
\section{Related Works}

\subsection{Foundational Financial Theories}

Traditional portfolio optimization methods, such as Markowitz's Modern Portfolio Theory (MPT) \citep{markowitz1952portfolio}, have been the industry standard for decades to balance risk and return through diversification. While effective in theory, MPT relies on static rules that often struggle to keep up with the chaotic, fast-moving reality of today's markets.

With the rise of machine learning, new approaches have emerged to tackle these limitations. \citet{jang2023deep} proposed integrating Deep Reinforcement Learning (DRL) with Technical Analysis and Modern Portfolio Theory constraints. Unlike traditional supervised learning models that merely predict prices, DRL agents learn sequential decision-making policies directly through interaction with the market environment. This transition from static allocation to dynamic portfolio management is emphasized in the comprehensive survey by \citet{pippas2025evolution}, which traces the evolution of reinforcement learning techniques in quantitative finance.

\subsection{Sequential Modeling and Temporal Dependencies}

Addressing the partial observability of financial markets, \citet{li2015recurrent} introduced a hybrid recurrent reinforcement learning framework integrating RNN-LSTM architectures within a Q-learning setup. Their approach enables agents to maintain memory of historical market states, capturing temporal dependencies that are critical in stochastic and non-stationary environments.

\subsection{The Evolution of Deep Reinforcement Learning in Finance}

Building upon classical frameworks, \citet{jang2023deep} combined DRL with MPT constraints using tensor decomposition techniques to manage high-dimensional financial data within structured risk frameworks. Similarly, \citet{jiang2017deep} proposed a CNN-based reinforcement learning framework that automates feature extraction from multi-channel price representations, reducing reliance on handcrafted indicators and improving scalability.

\subsection{Actor-Critic Frameworks}

Modern portfolio optimization strategies predominantly rely on the Actor-Critic architecture, which simultaneously updates a policy network (Actor) and a value function estimator (Critic). As implemented in the FinRL framework \citep{liu2020finrl}, in the ensemble strategy proposed by \citet{yang2025ensemble}, and discussed in \citet{jang2023deep}, actor-critic methods provide improved stability and continuous control capabilities.

\begin{itemize}
    \item \textbf{DDPG} (Deep Deterministic Policy Gradient) enables continuous portfolio weight allocation and deterministic policy updates.
    \item \textbf{A2C} (Advantage Actor-Critic) reduces gradient variance using an advantage function, improving robustness in volatile markets.
    \item \textbf{PPO} (Proximal Policy Optimization) stabilizes training by constraining policy updates through clipped objective functions.
\end{itemize}

\subsection{Alternative Data and Future Directions}

Recent work incorporates alternative data streams to enrich state representations. \citet{nawathe2024multimodal} integrated financial news sentiment using NLP models such as FinBERT into DRL-based portfolio optimization. Emerging research by \citet{liu2025quantum} explores quantum-enhanced reinforcement learning combined with LSTM forecasting signals to improve convergence stability and performance in noisy financial systems.

\subsection{Risk-Aware Rewards and Responsible Management}

To ensure robust and responsible financial decision-making, recent literature emphasizes risk-aware reward functions. Both the FinRL framework \citep{liu2020finrl} and the ensemble strategy proposed by \citet{yang2025ensemble} incorporate market frictions such as transaction costs and financial turbulence indices. These mechanisms penalize excessive speculation and enable agents to survive extreme market conditions, including financial crashes.



\bibliographystyle{plainnat}
\bibliography{references}
\end{document}